% 中文节标题：引言
\section{Introduction}

% 近年来，流式四维重建方法已经能够在顺序处理单目视频的过程中，联合恢复人体运动、相机轨迹和场景几何。这些方法将历史观测压缩到递归状态中，从而实现高效的在线推理，但其建模通常假设相邻帧来自一次连续拍摄。
% 经过剪辑的视频经常违背这一假设：它们由时间上依次出现的多个镜头组成，镜头之间的切换可能突然改变视角、消除视觉重叠，并显著改变可见场景与人体外观。因此，在边界之前积累的历史可能不再适合作为新镜头的递归上下文。
% 本文研究这种非连续输入下的流式重建。这里的“多镜头”是指单目视频中依次出现的镜头片段，而不是同一时刻的同步多视角观测。我们的目标是将每个镜头中的相机轨迹、重建人体和场景几何表示在一个统一世界坐标系中；该坐标系由第一镜头初始化，并作为整段序列共享的重建参考系。
Recent advances in streaming 4D reconstruction have made it possible to
jointly recover human motion, camera trajectories, and scene geometry from a
monocular video as it is processed sequentially
~\citep{wang2025cut3r,chen2026human3r,sur2026unicon3r}. These methods achieve
efficient online inference by compressing past observations into a recurrent
state, but their formulation largely assumes that adjacent frames belong to a
continuous capture. Edited videos routinely violate this assumption: they
comprise temporally successive shots separated by transitions that may
abruptly change viewpoint, remove visual overlap, and alter the visible scene
and human appearance
~\citep{pavlakos2022multishot,zhang2025humanmm,on2026multithumbs}. The history
accumulated before such a boundary may therefore no longer provide valid
recurrent context for the new shot. We study streaming reconstruction under
these discontinuities. Here, \emph{multi-shot} refers to successive segments
of a monocular video, rather than synchronized multi-view observations. The
goal is to express the camera trajectories, reconstructed humans, and scene
geometry of every shot in a \emph{common world frame}---the reconstruction
coordinate system initialized by the first shot and shared across the
sequence.

% 现有方法只解决了这一问题的一个方面。多镜头人体重建方法通过联合处理多个已有镜头或进行序列级优化，在镜头切换前后建立人物关系。这些方法能够有效恢复跨镜头结构，但通常并非为仅依靠有限递归状态、且不访问未来帧的逐步预测而设计。
% 有状态在线重建方法提供了与之互补的能力：它们按照时间顺序处理输入，并避免逐视频的全局优化。然而，其递归状态被设计为在连续观测中逐步演化，并没有显式处理历史有效性在镜头边界处突然改变的情况。
% 因此，这两类方法都没有直接满足本文的核心需求：在保留先前镜头所建立空间参照的同时，避免将不兼容的历史继续传播到新镜头的重建过程中。
Existing methods address only one side of this setting. Multi-shot human
reconstruction approaches relate people across transitions by reasoning over
multiple available shots, often through joint processing or sequence-level
optimization~\citep{pavlakos2022multishot,zhang2025humanmm,on2026multithumbs}.
Although effective for recovering inter-shot structure, they are generally
not designed to produce each estimate from a bounded recurrent state without
future-frame access. Stateful online reconstruction methods provide the
complementary capability: they process frames in temporal order and avoid
per-video global optimization
~\citep{wang2025cut3r,chen2026human3r,chen2025long3r}. Their recurrent
state, however, is designed to evolve through continuous observations and
provides no explicit mechanism for a sudden change in history validity at a
shot boundary. Neither line of work therefore addresses the central
requirement here: preserving the spatial reference established by earlier
shots while preventing incompatible history from being propagated through
reconstruction of the new shot.

% 这一缺口源于递归历史在镜头边界处的使用冲突。跨越边界继续传播状态，可以保留先前观测所建立的坐标参照，却也迫使新镜头与在非连续视角下积累的历史发生交互；切换处引入的误差因而可能在后续递归更新中持续传播。
% 重新初始化状态可以消除这种干扰，使新镜头内部恢复连贯的递归重建，但同时也会抹去将该镜头定位到统一世界坐标系所需的参照。因此，关键问题并不是应该保留还是丢弃历史，而是历史的哪些作用应当跨越镜头边界继续存在。
% 我们的核心观察是，历史只应作为临时的对齐参照继续发挥作用：它可以被读取以建立镜头间关系，但不应作为新镜头的递归状态继续传播。使用相同重建模型进行的受控比较支持这一判断：状态延续造成的是随时间变化的相机漂移，而不只是一次固定偏移；独立重建则能够保持新镜头内部的连贯性，却失去了与先前世界坐标系的空间关系。
This gap arises from a conflict in how recurrent history is used at a shot
boundary. Carrying the state across the boundary preserves the coordinate
reference established by previous observations, but forces the new shot to
interact with history accumulated under a discontinuous viewpoint; errors
introduced at the transition can then persist through subsequent recurrent
updates. Reinitializing the state removes this interference and restores
coherent recurrence within the new shot, but also erases the reference needed
to locate that shot in the common world frame. The key question is therefore
not whether history should be retained or discarded, but which of its roles
should survive the boundary. Our central insight is that history should
survive only as a temporary alignment reference: it may be read to establish
the inter-shot relation, but should not be propagated as the recurrent state
of the new shot. A controlled comparison using the same reconstruction model
supports this distinction: state carry-over produces time-varying camera
drift rather than a single fixed offset, whereas independent reconstruction
remains coherent within the new shot but loses its relation to the preceding
world frame (Section~\ref{sec:ablation}).

% 为解决这一冲突，我们提出 Shot3R：一个为递归历史的两种作用赋予不同信息生命周期的流式重建框架。在每个镜头边界处，Shot3R 临时读取先前状态以推断镜头间的空间关系，但不会将这一过程产生的状态保留为后续重建所使用的记忆。
% 新镜头内部的递归过程从独立初始化的状态开始，并且只有该状态会继续向后传播。在建立空间关系之后，基于几何的人物关联维持镜头边界两侧的身份连续性，相机—人体共享变换则将新镜头的重建结果表示到统一世界坐标系中。
% Shot3R 不需要访问未来帧、修改已经输出的估计结果或执行逐视频优化。在三个多人体基准上，相比连续传播历史状态，Shot3R 将按数据集等权计算的平均 W-MPJPE 降低了 9.5\%，并一致改善了相机精度和身份连续性；在 EgoBody 预定义的极端视角子集上，W-MPJPE 的降幅达到 29.9\%。
To resolve this conflict, we introduce \method{}, a streaming reconstruction
framework that assigns different lifetimes to the two roles of recurrent
history. At each shot boundary, \method{} temporarily consults the preceding
state to infer the inter-shot relation, without retaining the resulting state
as memory for subsequent reconstruction. Recurrence within the new shot
instead starts from an independently initialized state, which becomes the only
state propagated forward. Once the spatial relation is established,
geometry-based person association maintains identities across the boundary,
and a shared camera--human transform expresses the new-shot reconstruction in
the common world frame. \method{} operates without future-frame access,
revision of previously emitted estimates, or per-video optimization. Across
three multi-person benchmarks, it reduces the dataset-balanced mean W-MPJPE by
9.5\% over continuous state propagation and consistently improves camera
accuracy and identity continuity; on EgoBody's predefined extreme-viewpoint
subset, the W-MPJPE reduction reaches 29.9\%.

% 本文的主要贡献如下。
Our contributions are:
\begin{itemize}
  % 我们识别并刻画了递归历史在镜头边界处的使用冲突：历史是建立相邻镜头关系所必需的参照，但将其继续传播到新镜头可能引入持续的重建误差，而将其完全丢弃又会失去镜头间的空间参照。
  \item We identify and characterize a conflict in recurrent history at shot
  boundaries: history is necessary for relating successive shots, yet
  propagating it into the new shot can introduce persistent reconstruction
  errors, while discarding it entirely removes the inter-shot spatial
  reference.

  % 我们提出 Shot3R，一个将“读取历史以进行对齐”与“传播状态以完成重建”分开处理的流式多镜头重建框架。Shot3R 仅将边界之前的历史保留为临时参照，并通过独立初始化的递归状态推进新镜头的重建，同时维持镜头边界两侧的空间与身份一致性。
  \item We introduce \method{}, a streaming multi-shot reconstruction
  framework that separates how history is read for alignment from how state is
  propagated for reconstruction. \method{} retains pre-boundary history only
  as a temporary reference and advances the new shot using independently
  initialized recurrence, while maintaining spatial and identity consistency
  across the boundary.

  % 通过使用相同重建模型的受控比较以及三个多人体基准上的评测，我们证明这种分离能够改善世界坐标人体重建、相机精度和身份连续性，并在大幅视角变化下取得最明显的收益，同时不需要访问未来帧或执行逐视频优化。
  \item Through controlled same-model comparisons and evaluations on three
  multi-person benchmarks, we show that this separation improves world-frame
  human reconstruction, camera accuracy, and identity continuity, with its
  largest gains under substantial viewpoint changes and without future-frame
  access or per-video optimization.
\end{itemize}
